\documentclass[a4paper, 11pt]{ltjsreport}
\usepackage{luatexja}
\usepackage{amssymb, amsmath}   % 数理環境用パッケージ
\usepackage{amsthm}             % 定理環境用パッケージ
\usepackage{ascmac}             % box環境用パッケージ
% \usepackage[dvipdfmx]{graphicx}           % 画像環境用パッケージ
\usepackage{mathtools}
\usepackage{amsfonts}
\usepackage{enumerate}
\usepackage{bm}
\usepackage{cite}%参考文献
\usepackage{physics}%ノルムとかが書きやすくなる。
\usepackage{listings}
\usepackage{jvlisting}
\theoremstyle{plain}
\theoremstyle{definition}
\newtheorem{公理}{公理}
\newtheorem{定義}{定義}
\newtheorem{定理}{定理}
\newtheorem{証明}{証明}
\newtheorem*{メモ}{メモ}

\newtheorem{dfn}{Definition}

\begin{document}
\chapter{現状}
\begin{equation}
  \left\{ \,
      \begin{aligned}
          -\Delta u &= \frac{1}{\epsilon^2}\qty(-\alpha u +\beta u^2 -\gamma u^3 -\epsilon\qty(a v + b w ))\\
          -\Delta v &= u-v \\
          -\Delta w &= \frac{1}{d^2}\qty(u-w)\\
          \norm{u}_{L^2} &= N
          %\dot{v}&=v-\dfrac{v^3}{3}-w+I_{ext}\\
          %\tau \dot{w}&=v-a-bw
      \end{aligned}
  \right.
\end{equation}
\begin{equation}
  \left\{ \,
      \begin{aligned}
          -\Delta u &= \frac{1}{\epsilon^2}\qty(-\alpha u +\beta u^2 -\gamma u^3 -\epsilon\qty(a v + b w ))\\
          -\Delta v &= u-v \\
          -\Delta w &= \frac{1}{d^2}\qty(u-w)\\
          \sqrt{\qty( u, u)_{L^2} } &= N
          %\dot{v}&=v-\dfrac{v^3}{3}-w+I_{ext}\\
          %\tau \dot{w}&=v-a-bw
      \end{aligned}
  \right.
\end{equation}
\begin{align}
   \bm{F}(\bm{u},\epsilon) =Ku_h- \dfrac{-\alpha u_h + \beta  u_h^2 -\gamma  u_h^3-\epsilon(aL(K+L)^{-1}L+bL(d^2K+L)^{-1}) u_h}{\epsilon^2}
\end{align}
\begin{equation}
    \bm{\mathcal{F}}(\bm{X}) = 
    \begin{pmatrix}
        \bm{F}(\bm{u}, \epsilon) \\
        \sqrt{\bm{u}^T L \bm{u}} - N
    \end{pmatrix}
    = \bm{0}
\end{equation}

\newpage
\begin{equation}
    \begin{pmatrix}
        \dfrac{\partial \bm{F}}{\partial \bm{u}} & \dfrac{\partial \bm{F}}{\partial \epsilon} \\[10pt]
        \dfrac{\partial g}{\partial \bm{u}} & \dfrac{\partial g}{\partial \epsilon}
    \end{pmatrix}
    \begin{pmatrix}
        \Delta \bm{u} \\
        \Delta \epsilon
    \end{pmatrix}
    =
    \begin{pmatrix}
        -\bm{F}(\bm{u}, \epsilon) \\
        -g(\bm{u})
    \end{pmatrix}
\end{equation}

\begin{align}
    \dfrac{\partial \bm{F}}{\partial \bm{u}} 
    = 
    % K - \dfrac{1}{\epsilon^2} \left( -\alpha L + 2\beta L u_h - 3\gamma L u_h^2 - \delta L(K+rL)^{-1}L \right)
    K- \dfrac{-\alpha L + 2\beta  L u_h -3\gamma L u_h^2-\epsilon(aL(K+L)^{-1}L+bL(d^2K+L)^{-1}L) }{\epsilon^2}
\end{align}
\begin{align}
  \dfrac{\partial \bm{F}}{\partial \epsilon}
  =
  \dfrac{2(-\alpha L u_h + \beta  u_h^2 -\gamma  u_h^3)}{\epsilon^3}
  -\dfrac{(aL(K+L)^{-1}L+bL(d^2K+L)^{-1}L)u_h}{\epsilon^2}
\end{align}
\begin{align}
  \dfrac{\partial g}{\partial \bm{u}}= \dfrac{u^T L}{2\sqrt{u^T L u}}
\end{align}
\begin{align}
  \dfrac{\partial g}{\partial \epsilon}=0
\end{align}
\end{document}